%-------------------------
% Resume in Latex
% Author : Jake Gutierrez
% Based off of: https://github.com/sb2nov/resume
% License : MIT
%------------------------

\documentclass[letterpaper,11pt]{article}

% Keep the original serif appearance while encoding accented characters directly.
\usepackage[T1]{fontenc}
\usepackage[utf8]{inputenc}
\usepackage{lmodern}
\usepackage{microtype}
\usepackage{latexsym}
\usepackage[empty]{fullpage}
\usepackage{titlesec}
\usepackage{marvosym}
\usepackage[dvipsnames]{xcolor}
\usepackage{verbatim}
\usepackage{enumitem}
\usepackage{hyperref}
\usepackage{fancyhdr}
\usepackage[english]{babel}
\input{glyphtounicode}

% Two coordinated accent colors; body text and publication titles remain black.
% Edit these HEX values to change the palette.
\definecolor{HeadingNavy}{HTML}{334E68}
\definecolor{AuthorTeal}{HTML}{14665E}
\hypersetup{
  colorlinks=true,
  urlcolor=HeadingNavy,
  linkcolor=HeadingNavy,
  pdftitle={Luis F. Recalde - Curriculum Vitae},
  pdfauthor={Luis F. Recalde}
}
\newcommand{\authorname}{\textcolor{AuthorTeal}{\textbf{Luis F. Recalde}}}
\newcommand{\journal}[1]{\textcolor{HeadingNavy}{#1}}
\newcommand{\papertitle}[1]{\textcolor{black}{\textbf{#1}}}
% Shared destination for the name and Website links in the header.
\newcommand{\websiteURL}{https://lfrecalde1.github.io/}


%----------FONT OPTIONS----------
% sans-serif
% \usepackage[sfdefault]{FiraSans}
% \usepackage[sfdefault]{roboto}
% \usepackage[sfdefault]{noto-sans}
% \usepackage[default]{sourcesanspro}

% serif
% \usepackage{CormorantGaramond}
% \usepackage{charter}


\setlength{\footskip}{18pt}
\pagestyle{fancy}
\fancyhf{} % clear all header and footer fields
\fancyfoot[R]{\footnotesize\thepage}
\renewcommand{\headrulewidth}{0pt}
\renewcommand{\footrulewidth}{0pt}

% Adjust margins
\addtolength{\oddsidemargin}{-0.5in}
\addtolength{\evensidemargin}{-0.5in}
\addtolength{\textwidth}{1in}
\addtolength{\topmargin}{-.5in}
\addtolength{\textheight}{1.0in}
\setlength{\headwidth}{\textwidth}

\urlstyle{same}

\raggedbottom
% Use full justification for body paragraphs without paragraph indentation.
\setlength{\parindent}{0pt}

% Light section bands separate the main sections without adding more accent colors.
\newcommand{\sectionband}[1]{%
  \begingroup
  \setlength{\fboxsep}{3pt}%
  \colorbox{HeadingNavy!6}{%
    \makebox[\dimexpr\linewidth-2\fboxsep\relax][l]{%
      \textcolor{AuthorTeal}{\rule[-2pt]{2pt}{12pt}}\hspace{6pt}%
      \MakeUppercase{#1}}}%
  \endgroup
}
\titleformat{\section}{
  \sffamily\bfseries\upshape\raggedright\fontsize{13}{15}\selectfont\color{HeadingNavy}
}{}{0em}{\sectionband}
\titlespacing*{\section}{0pt}{12pt}{5pt}

% Preserve Unicode text mappings for PDF text extraction (not an ATS guarantee).
\pdfgentounicode=1
% Emit explicit word spaces as well as positioning them visually.
\pdfinterwordspaceon

%-------------------------
% Custom commands
\newcommand{\resumeItem}[1]{
  \item {\small #1\par}\vspace{-2pt}
}

% Ordinary paragraphs keep each date beside its role or institution in extracted text.
\newcommand{\resumeSubheading}[4]{
  \vspace{-2pt}\item
    \textbf{#1}\enspace\textbar\enspace #2\\
    {\small\textit{#3\enspace\textbar\enspace #4}\par}\vspace{-5pt}
}

\newcommand{\resumeSubSubheading}[2]{
    \item {\small\textit{#1\enspace\textbar\enspace #2}\par}\vspace{-5pt}
}

\newcommand{\resumeProjectHeading}[2]{
    \item {\small #1\enspace\textbar\enspace #2\par}\vspace{-5pt}
}

\newcommand{\resumeSubItem}[1]{\resumeItem{#1}\vspace{-4pt}}

\renewcommand\labelitemi{\textbullet}
\renewcommand\labelitemii{\textbullet}

\newcommand{\resumeSubHeadingListStart}{\begin{itemize}[leftmargin=0.15in, label={}, topsep=3pt]}
\newcommand{\resumeSubHeadingListEnd}{\end{itemize}}
\newcommand{\resumeItemListStart}{\begin{itemize}}
\newcommand{\resumeItemListEnd}{\end{itemize}\vspace{-5pt}}

%-------------------------------------------
%%%%%%  RESUME STARTS HERE  %%%%%%%%%%%%%%%%%%%%%%%%%%%%


\begin{document}
\begin{center}
\href{\websiteURL}{\textcolor{AuthorTeal}{\textbf{\Huge Luis F. Recalde}}} \\ \vspace{3pt}
\small Robotics \& Controls Research Engineer $|$ Nonlinear Control, Optimization \& Autonomous Navigation \\ \vspace{3pt}
Worcester, MA $|$ \href{mailto:lfrecalde@wpi.edu}{lfrecalde@wpi.edu} $|$
\href{https://github.com/lfrecalde1}{GitHub} $|$
\href{https://www.linkedin.com/in/luis-f-recalde-8a6432123}{LinkedIn} $|$
\href{\websiteURL}{Website} $|$
\href{https://scholar.google.com/citations?user=_v1511cAAAAJ}{Google Scholar}
\end{center}
\vspace{-12pt}
\section{Research Focus}
{\small I am a research engineer and PhD candidate in Robotics Engineering at Worcester Polytechnic Institute.
I hold an MSc in Control Systems Engineering and a bachelor's degree in Mechatronics Engineering.
My research aims to enable \textbf{fully onboard autonomous navigation in unstructured environments}.
I develop control and estimation algorithms using \textbf{nonlinear control theory, differential geometry, and numerical optimization}.
I take these methods from dynamical models and physics-based simulation to \textbf{real-world experiments}.\par}
\section{Research Experience}
\resumeSubHeadingListStart
\resumeSubheading{Research Assistant}{Jan 2025 -- Present}
{Worcester Polytechnic Institute $|$ \href{https://acp-lab.github.io/acplab-web/}{Aerial-Robot Control and Perception Lab}}{Worcester, MA}
\resumeItemListStart
\resumeItem{Developed \href{https://acp-lab.github.io/acplab-web/}{ACP Lab}'s \textbf{C++/ROS 2 MuJoCo simulator} for aerial vehicles and cable-suspended payloads, combining aerodynamic and motor dynamics with simulated cameras and LiDAR.}
\resumeItem{Implemented \textbf{incremental nonlinear dynamic inversion (INDI)} body-rate control and a communication protocol in Betaflight.}
\resumeItemListEnd
\resumeSubheading{Visiting Researcher}{May 2024 -- Dec 2024}
{New York University $|$ \href{https://arplaboratory.github.io/}{Agile Robotics and Perception Lab}}{New York, NY}
\resumeItemListStart
\resumeItem{Developed controllers using \textbf{control Lyapunov functions (CLFs)} and \textbf{control barrier functions (CBFs)} for quadrotors with cable-suspended payloads, proving \textbf{stability and camera field-of-view safety guarantees}.}
\resumeItemListEnd
\resumeSubheading{Research Engineer}{Mar 2022 -- Mar 2023}
{Universidad Indoam\'erica $|$ SISAU Research Group and CICHE}{Ecuador}
\resumeItemListStart
\resumeItem{Applied optimization methods to robotic systems.}
\resumeItemListEnd
\resumeSubheading{Research Engineer}{May 2021 -- Oct 2021}
{Escuela Superior Polit\'ecnica de Chimborazo}{Ecuador}
\resumeItemListStart
\resumeItem{Applied optimization methods to networked Internet of Things (IoT) sensing for road-safety applications.}
\resumeItemListEnd
\resumeSubHeadingListEnd
\section{Technical Skills}
{\small\textbf{Optimization}\par}
\begin{itemize}[leftmargin=0.22in, itemsep=2pt, parsep=0pt, topsep=4pt]
\small
\item \textbf{Model predictive control (MPC):} \href{https://arxiv.org/abs/2510.11525}{Nonlinear model predictive control (NMPC)} and \href{https://arxiv.org/abs/2609.36482}{model predictive contouring control (MPCC)}; solution methods include sequential quadratic programming (SQP), SQP with real-time iteration (SQP-RTI), interior-point methods, differential dynamic programming (DDP), and iterative linear quadratic regulator (iLQR).
\item \textbf{Nonsmooth and hybrid optimal control:} complementarity-constrained formulations and relaxation methods for trajectory optimization across hybrid mode transitions.
\item \textbf{Manifold optimization:} Geometric formulations for \href{https://arxiv.org/abs/2410.13541}{robot pose estimation} and \href{https://arxiv.org/abs/2510.11525}{control}.
\item \textbf{Differentiable optimization:} Implicit Karush--Kuhn--Tucker (KKT) differentiation to compute solution sensitivities for gradient-based MPC tuning and learning.
\item \textbf{Distributed estimation and control:} Alternating direction method of multipliers (ADMM) for distributed sensor fusion and multi-robot coordination.
\end{itemize}
{\small\textbf{Modeling and Control}\par}
\begin{itemize}[leftmargin=0.22in, itemsep=2pt, parsep=0pt, topsep=4pt]
\small
\item \textbf{Geometric control on Lie groups:} Controller design and \href{https://arxiv.org/abs/2504.08841}{stability analysis} for \href{https://arxiv.org/abs/2510.11525}{attitude and pose tracking}.
\item \textbf{Task-priority control:} Null-space projection and affine-subspace methods for prioritizing robotic tasks (\href{https://arxiv.org/abs/2311.08019}{unmanned aerial vehicle (UAV) inspection}; \href{https://arxiv.org/abs/2605.20392}{tactile manipulation}).
\item \textbf{Safety-critical control:} \href{https://arxiv.org/abs/2504.08841}{CLFs} and \href{https://arxiv.org/abs/2504.08841}{CBFs} for robotic systems.
\item \textbf{Data-driven modeling and control:} \href{https://www.mdpi.com/1424-8220/22/13/4712}{Dynamic mode decomposition with control (DMDc)} and Koopman operator approximations for robot dynamics identification and predictive modeling.
\end{itemize}
\newpage
\section{Technical Skills}
{\small\textbf{Perception and State Estimation}\par}
\begin{itemize}[leftmargin=0.22in, itemsep=2pt, parsep=0pt, topsep=4pt]
\small
\item \textbf{Visual and tactile servoing:} Image-based feedback with NMPC for \href{https://doi.org/10.1109/LRA.2024.3360876}{UAV photovoltaic array inspection}; vision-based tactile features and MPC for \href{https://doi.org/10.1109/LRA.2026.3698362}{manipulator contour following}.
\item \textbf{Odometry and sensor fusion:} \href{https://doi.org/10.1016/j.robot.2025.105009}{Dual-quaternion LiDAR odometry and mapping}; contracting nonlinear observers for LiDAR--inertial fusion.
\end{itemize}

\section{Education}
\resumeSubHeadingListStart
\resumeSubheading{Worcester Polytechnic Institute}{Jan 2025 -- Present}
{PhD in Robotics Engineering, in progress $|$ Advisor: Prof. Guanrui Li}{Worcester, MA}
\resumeSubheading{Universidad Nacional de San Juan (INAUT)}{May 2023 -- Dec 2024}
{MSc in Control Systems Engineering $|$ DAAD scholarship recipient}{Argentina}
\resumeSubheading{Universidad de las Fuerzas Armadas--ESPE}{2021}
{Bachelor's Degree in Mechatronics Engineering}{Ecuador}
\resumeSubHeadingListEnd

\section{Software and Tools}
\begin{itemize}[leftmargin=0.15in, label={}, itemsep=2pt, parsep=0pt, topsep=3pt]
\small
\item \textbf{Programming:} C++, Python, MATLAB.
\item \textbf{Simulation:} MuJoCo (C++ integration and plugins), MJX (JAX backend), Newton, Webots.
\item \textbf{Optimization:} CasADi, acados, Crocoddyl, IPOPT, HPIPM, OSQP, LCQPow.
\item \textbf{Middleware and Deployment:} Robot Operating System 2 (ROS 2), Zenoh, Linux, Git, Docker, Singularity.
\item \textbf{Flight Systems:} Betaflight, PX4.
\end{itemize}

\section{Publications}
\begin{enumerate}[leftmargin=0.28in, label={[\arabic*]}, itemsep=3pt, parsep=0pt, topsep=3pt]
\small
\item \href{https://arxiv.org/abs/2609.36482}{\papertitle{DQ-MPCC: Dual-Quaternion MPCC for Quadrotor Racing.}}\\
B. S. Guevara, \authorname, G. Li, T. Nascimento. \journal{IEEE RA-L}, accepted, 2026.
\item \href{https://doi.org/10.1109/LRA.2026.3698362}{\papertitle{VBT-MPC: Vision-Based Tactile MPC for Contour Following.}}\\
E. Velasco-S\'anchez, \authorname, G. Li, P. Gil. \journal{IEEE RA-L}, 11(7), 8504--8511, 2026.
\item \href{https://doi.org/10.1109/LRA.2025.3632070}{\papertitle{ES-HPC-MPC: Exponentially Stable Hybrid Perception Constrained MPC for Quadrotor with Suspended Payloads.}}\\
\authorname, M. Sarvaiya, G. Loianno, G. Li. \journal{IEEE RA-L}, 11(1), 266--273, 2026.
\item \href{https://doi.org/10.1109/LRA.2025.3623022}{\papertitle{DQ-NMPC: Dual-Quaternion NMPC for Quadrotor Flight.}}\\
\authorname, D. Agrawal, J. Arrizabalaga, G. Li. \journal{IEEE RA-L}, 2025.
\item \href{https://doi.org/10.1016/j.robot.2025.105009}{\papertitle{DualQuat-LOAM: LiDAR Odometry and Mapping Parametrized on Dual Quaternions.}}\\
E. P. Velasco-S\'anchez, \authorname\ et al. \journal{Robotics and Autonomous Systems}, 191, 105009, 2025.
\item \href{https://doi.org/10.1109/LRA.2024.3360876}{\papertitle{Visual Servoing NMPC Applied to UAVs for Photovoltaic Array Inspection.}}\\
E. P. Velasco-S\'anchez, \authorname\ et al. \journal{IEEE RA-L}, 9(3), 2766--2773, 2024.
\item \href{https://doi.org/10.3390/drones7020144}{\papertitle{A Comparative Study between NMPC and Baseline Feedback Controllers for UAV Trajectory Tracking.}}\\
B. S. Guevara, \authorname\ et al. \journal{Drones}, 7(2), 144, 2023.
\item \href{https://doi.org/10.3390/s22134712}{\papertitle{System Identification and Nonlinear Model Predictive Control with Collision Avoidance Applied in Hexacopters UAVs.}}\\
\authorname\ et al. \journal{Sensors}, 22(13), 4712, 2022.
\item \href{https://doi.org/10.3390/s21093057}{\papertitle{Three-Dimensional Unified Motion Control of a Robotic Standing Wheelchair for Rehabilitation Purposes.}}\\
J. S. Ortiz, G. Palacios-Navarro, V. H. Andaluz, \authorname. \journal{Sensors}, 21(9), 3057, 2021.
\end{enumerate}
\section{Service}
{\small
\textbf{Peer Review:} IEEE Robotics and Automation Letters (RA-L); IEEE Control Systems Letters (L-CSS); Robotics and Autonomous Systems (RAS);\\
IEEE International Conference on Robotics and Automation (ICRA); IEEE/RSJ International Conference on Intelligent Robots and Systems (IROS).\par}
\end{document}
